\documentclass[11pt]{article}

\usepackage{amsmath,amssymb,amsthm}

\newtheorem{lemma}{Lemma}
\newtheorem{theorem}{Theorem}

\title{Tube Containment Reduction for Cyclone}
\date{}

\begin{document}

\maketitle

\section*{Setup}

Let \(G\) be a graph of maximum degree \(\Delta\).  
Let \(C_1,C_2\) be simple cycles with identical local signatures under
\(\mathrm{FO}^k\) radius-\(R\) neighborhood types.  
Let
\[
D = C_1 \oplus C_2 .
\]
Set
\[
r=\min(R/2,3^k).
\]

\begin{lemma}[Tube Containment Reduction]
Assume every connected component of \(D\) is contained in a union of at most
\(m=m(k,\Delta,R)\) radius-\(r\) balls centered on vertices of \(C_1\cup C_2\).
Then every connected component of \(D\) has size bounded by
\[
L=L(k,\Delta,R):=m \cdot \max_{v\in V(G)} |B_r(v)|.
\]
\end{lemma}

\begin{proof}
Each radius-\(r\) ball contains at most
\[
1+\Delta+\Delta^2+\cdots+\Delta^r
\]
vertices. Hence any component contained in a union of at most \(m\) such balls
has at most
\[
m(1+\Delta+\Delta^2+\cdots+\Delta^r)
\]
vertices. Since every connected component of \(D\) is 2-regular, each component
is a cycle, so its length is bounded by the same quantity.
\end{proof}

\begin{theorem}[Reduction to Tube Number Bound]
If there exists \(m=m(k,\Delta,R)\) such that every connected component of
\(C_1\oplus C_2\) is contained in a union of at most \(m\) radius-\(r\) balls,
then the Short-Cycle Symmetric-Difference Lemma holds.
\end{theorem}

\begin{proof}
Apply the Tube Containment Reduction lemma.
\end{proof}

\end{document}
